\documentclass{article}
\usepackage[utf8]{inputenc}
\usepackage{geometry}
\geometry{a4paper, margin=1in}

\title{Análisis del Proyecto de Seguridad Informática}
\author{Nathaniel}
\date{\today}

\begin{document}

\maketitle

\section{Resumen del Aporte}
El aporte de Nathaniel al proyecto consistió en desarrollar la funcionalidad para consultar registros DNS de tipo MX y NS. Implementó la función \texttt{get\_mx\_ns\_records()} en el archivo \texttt{modulos.py} y la integró con el script principal \texttt{auditorias.py} a través del argumento de línea de comandos \texttt{--dns-mxns}.

\section{Análisis del Código}

\subsection{Puntos Fuertes}
\begin{itemize}
    \item \textbf{Funcionalidad Correcta:} La tarea asignada fue completada y la funcionalidad se integra correctamente en el sistema general.
    \item \textbf{Manejo de Errores:} La función incluye un buen manejo de excepciones para los casos en que no se encuentran los registros DNS solicitados, lo que hace al código más robusto.
    \item \textbf{Claridad:} El código es legible y está bien estructurado.
\end{itemize}

\subsection{Áreas de Mejora (Diseño de Software)}
\begin{itemize}
    \item \textbf{Separación de Lógica y Presentación:} La función \texttt{get\_mx\_ns\_records()} actualmente imprime los resultados directamente en la consola. Una mejor práctica de diseño de software es que las funciones de lógica retornen los datos (por ejemplo, en una lista o un diccionario). El script que llama a la función (\texttt{auditorias.py}) debería ser el responsable de presentar esos datos al usuario. Esto aumenta la modularidad y permite que la función se reutilice más fácilmente en otros contextos.
\end{itemize}

\section{Conclusión}
El trabajo de Nathaniel es de buena calidad y cumple con los requisitos de la tarea. La sugerencia de mejora se enfoca en un principio de diseño de software más avanzado que puede fortalecer aún más la estructura del código.

\end{document}
